% ===================================================================
%  MTH2021 Linear Algebra with Applications  —  Shared Preamble
% ===================================================================
\usepackage[a4paper,margin=1.5cm,headheight=14pt]{geometry}
\usepackage{amsmath,amssymb,amsthm,mathtools}
\usepackage{enumitem}
\usepackage{xcolor}
\usepackage[most]{tcolorbox}
\usepackage{titlesec}
\usepackage[Conny]{fncychap}
\usepackage{fancyhdr}
\usepackage{graphicx}
\usepackage{adjustbox}
\usepackage{array}
\usepackage{booktabs}
\usepackage{multicol}
\usepackage[T1]{fontenc}
\usepackage[protrusion=true,expansion=false]{microtype}
\usepackage{hyperref}
\hypersetup{colorlinks=true,linkcolor=NavyBlue,citecolor=NavyBlue,
            urlcolor=NavyBlue,linktoc=all}

% ---- Palette ------------------------------------------------------
\definecolor{NavyBlue}{RGB}{16,52,96}
\definecolor{InkBlue}{RGB}{28,62,110}
\definecolor{BoxGrey}{RGB}{244,246,249}
\definecolor{RuleGrey}{RGB}{120,130,145}
\definecolor{ThmBlue}{RGB}{232,238,247}
\definecolor{DefGreen}{RGB}{234,243,236}
\definecolor{DefRule}{RGB}{52,120,72}
\definecolor{ExAmber}{RGB}{250,246,235}
\definecolor{ExRule}{RGB}{170,130,40}
\definecolor{WarnRed}{RGB}{251,239,239}
\definecolor{WarnRule}{RGB}{160,55,55}

% ---- Headers ------------------------------------------------------
\pagestyle{fancy}
\fancyhf{}
\renewcommand{\headrulewidth}{0.4pt}
\fancyhead[L]{\small\itshape\nouppercase{\leftmark}}
\fancyhead[R]{\small\thepage}
\fancyfoot[C]{\small MTH2021 Linear Algebra with Applications}

% ---- Section styling ---------------------------------------------
\titleformat{\section}{\Large\bfseries\color{NavyBlue}}{\thesection}{0.7em}{}
\titleformat{\subsection}{\large\bfseries\color{InkBlue}}{\thesubsection}{0.6em}{}
\titleformat{\subsubsection}{\normalsize\bfseries\color{InkBlue}}{\thesubsubsection}{0.5em}{}

% ===================================================================
%  Theorem-like environments — MANUAL numbering to mirror the
%  seminar notes exactly (e.g. Theorem 1.5.7, Definition 1.2.3).
%  Usage:  \begin{thms}{1.5.7}{Equivalent statements} ... \end{thms}
%  The second argument (title) may be empty: {...}{}
%  Auto-creates \label targets thm:1.5.7, def:1.2.3, etc.
% ===================================================================
\makeatletter
\newcommand{\@mlabeltitle}[1]{\if\relax\detokenize{#1}\relax\else\;(#1)\fi}

\newtcolorbox{resultbox}[3]{%
  boxrule=0pt, frame hidden, sharp corners,
  borderline west={3pt}{0pt}{InkBlue},
  colback=ThmBlue, breakable, enhanced jigsaw,
  left=10pt,right=8pt,top=6pt,bottom=6pt,before skip=10pt,after skip=10pt,
  before upper={\textbf{#1~#2}\@mlabeltitle{#3}\par\smallskip}}

\newenvironment{thms}[2]{%
  \hypertarget{thm:#1}{}\begin{resultbox}{Theorem}{#1}{#2}}{\end{resultbox}}
\newenvironment{lems}[2]{%
  \hypertarget{lem:#1}{}\begin{resultbox}{Lemma}{#1}{#2}}{\end{resultbox}}
\newenvironment{cors}[2]{%
  \hypertarget{cor:#1}{}\begin{resultbox}{Corollary}{#1}{#2}}{\end{resultbox}}
\newenvironment{props}[2]{%
  \hypertarget{prop:#1}{}\begin{resultbox}{Proposition}{#1}{#2}}{\end{resultbox}}

\newtcolorbox{defbox}[2]{%
  boxrule=0pt, frame hidden, sharp corners,
  borderline west={3pt}{0pt}{DefRule},
  colback=DefGreen, breakable, enhanced jigsaw,
  left=10pt,right=8pt,top=6pt,bottom=6pt,before skip=10pt,after skip=10pt,
  before upper={\textbf{Definition~#1}\@mlabeltitle{#2}\par\smallskip}}
\newenvironment{defns}[2]{%
  \hypertarget{def:#1}{}\begin{defbox}{#1}{#2}}{\end{defbox}}

\newtcolorbox{algbox}[2]{%
  boxrule=0pt, frame hidden, sharp corners,
  borderline west={3pt}{0pt}{RuleGrey},
  colback=BoxGrey, breakable, enhanced jigsaw,
  left=10pt,right=8pt,top=6pt,bottom=6pt,before skip=10pt,after skip=10pt,
  before upper={\textbf{Algorithm~#1}\@mlabeltitle{#2}\par\smallskip}}
\newenvironment{algs}[2]{%
  \hypertarget{alg:#1}{}\begin{algbox}{#1}{#2}}{\end{algbox}}

\newtcolorbox{rmkbox}[1]{%
  boxrule=0pt, frame hidden, sharp corners,
  borderline west={2pt}{0pt}{RuleGrey},
  colback=white, breakable, enhanced jigsaw,
  left=10pt,right=4pt,top=2pt,bottom=2pt,before skip=8pt,after skip=8pt,
  before upper={\textbf{Remark~#1.}\;}}
\newenvironment{rmks}[1]{%
  \hypertarget{rmk:#1}{}\begin{rmkbox}{#1}}{\end{rmkbox}}

% Unnumbered remark/aside (no "Remark N." prefix auto-added)
\newtcolorbox{rmknotebox}{%
  boxrule=0pt, frame hidden, sharp corners,
  borderline west={2pt}{0pt}{RuleGrey},
  colback=white, breakable, enhanced jigsaw,
  left=10pt,right=4pt,top=2pt,bottom=2pt,before skip=8pt,after skip=8pt}
\newenvironment{rmknote}{\begin{rmknotebox}}{\end{rmknotebox}}

\newtcolorbox{exmpbox}[2]{%
  colback=ExAmber, colframe=ExRule, boxrule=0.6pt, sharp corners,
  breakable, enhanced, left=9pt,right=9pt,top=7pt,bottom=7pt,
  before skip=11pt, after skip=11pt,
  fonttitle=\bfseries\color{white}, coltitle=white, colbacktitle=ExRule,
  attach boxed title to top left={yshift=-0.5pt,xshift=8pt},
  boxed title style={sharp corners,boxrule=0pt},
  title={Example~#1\@mlabeltitle{#2}}}
\newenvironment{exmp}[2]{%
  \hypertarget{eg:#1}{}\begin{exmpbox}{#1}{#2}}{\end{exmpbox}}
\makeatother

% Wide-display helper: shrink an over-wide math display to \linewidth.
\newcommand{\widedisplay}[1]{%
  \begin{center}\adjustbox{max width=\linewidth}{$\displaystyle #1$}\end{center}}

% Warning / pitfall box
\newtcolorbox{pitfall}[1][Warning]{
  colback=WarnRed, colframe=WarnRule, boxrule=0.6pt, sharp corners,
  breakable, enhanced, left=9pt,right=9pt,top=6pt,bottom=6pt,
  before skip=10pt, after skip=10pt,
  fonttitle=\bfseries, coltitle=white, colbacktitle=WarnRule, title={#1}}

% Chapter "roadmap" intro box
\newtcolorbox{roadmap}{
  colback=BoxGrey, colframe=NavyBlue, boxrule=0.8pt, sharp corners,
  breakable, enhanced, left=10pt,right=10pt,top=8pt,bottom=8pt,
  before skip=6pt, after skip=14pt}

% ---- Proof tweak --------------------------------------------------
\renewcommand{\qedsymbol}{$\blacksquare$}

% ===================================================================
%  Redesigned PROOF presentation.
%  proofbox : a light box with a left rule in the theorem blue, so it
%             reads as "hanging off" the result it proves. Optional
%             argument relabels the header (e.g. [Proof of (b)]).
%  steps    : one logical step per line; \st{statement}{reason} puts
%             the reason flush right (seminar "statement | reason").
%             \stn{statement} for a step with no cited reason.
%  case     : \case{(a)} heads a sub-part with space above it.
%  proofomit: deferral notice (MTH2025 / Problem Set), same styling.
% ===================================================================
\makeatletter
\definecolor{ProofBlue}{RGB}{60,96,150}
\definecolor{ProofBack}{RGB}{248,250,253}

\newtcolorbox{proofboxenv}[1]{%
  boxrule=0pt, frame hidden, sharp corners,
  borderline west={2.5pt}{0pt}{ProofBlue},
  colback=ProofBack, breakable, enhanced jigsaw,
  left=11pt,right=8pt,top=5pt,bottom=5pt,
  before skip=-4pt, after skip=12pt,
  before upper={\textit{\color{ProofBlue}#1.}\par\smallskip}}

% \begin{prf} ... \end{prf}            -> header "Proof"
% \begin{prf}[Proof of (b)] ... \end{prf}
\newenvironment{prf}[1][Proof]{\begin{proofboxenv}{#1}}{\unskip\nobreak\hfil\penalty50\hskip1em\null\nobreak\hfil$\blacksquare$\parfillskip=0pt\finalhyphendemerits=0\end{proofboxenv}}

% Step list: each step on its own line, reason flush right.
\newenvironment{steps}{%
  \begin{list}{}{%
    \setlength{\leftmargin}{1.2em}\setlength{\rightmargin}{0pt}%
    \setlength{\itemsep}{3pt}\setlength{\parsep}{0pt}%
    \setlength{\topsep}{3pt}\setlength{\labelwidth}{0pt}%
    \setlength{\labelsep}{0pt}\setlength{\listparindent}{0pt}}%
  }{\end{list}}
% \st{statement}{reason}
\newcommand{\st}[2]{\item $\displaystyle #1$ \hfill\makebox[0pt][r]{\small\color{RuleGrey}(#2)}\par}
% step with no reason
\newcommand{\stn}[1]{\item $\displaystyle #1$\par}
% inline text step (for prose lines inside a step list)
\newcommand{\stt}[1]{\item #1\par}

% Case header inside a proof. The label sits on its own, and the case body
% flows after it (possibly over several lines / steps) clearly grouped.
% \case{(a)} Body text...  -> "(a)  Body text..." with space above, hanging indent.
\newcommand{\case}[1]{%
  \par\addvspace{4pt}%
  \noindent\hangindent=1.6em\hangafter=1
  \textbf{\color{ProofBlue}#1}\hspace{0.6em}\ignorespaces}
% Use \endcase optionally before the next \case for clarity (no-op spacing helper).
\newcommand{\smallgap}{\par\addvspace{2pt}}

% Structured "check" list: for proofs that verify a list of named conditions
% (Subspace Test, independence checks, etc). Each item has a bold blue label
% followed by the verification on the same logical block.
%   \begin{checks}
%     \chk{Nonempty} ...justification...
%     \chk{Closure under addition} ...
%     \chk{Closure under scalar multiplication} ...
%   \end{checks}
% Follow with a concluding sentence ("Therefore ... is a subspace.").
\newenvironment{checks}{%
  \par\addvspace{3pt}%
  \begin{list}{}{%
    \setlength{\leftmargin}{1.5em}\setlength{\rightmargin}{0pt}%
    \setlength{\itemsep}{4pt}\setlength{\parsep}{2pt}%
    \setlength{\topsep}{2pt}\setlength{\labelwidth}{0pt}%
    \setlength{\labelsep}{0pt}\setlength{\listparindent}{0pt}}%
  }{\end{list}\par\addvspace{2pt}}
% \chk{Label} body...
\newcommand{\chk}[1]{\item \textbf{\color{ProofBlue}#1:}\enspace\ignorespaces}

% Deferral notice, styled like a faint proof box.
\newtcolorbox{omitbox}{%
  boxrule=0pt, frame hidden, sharp corners,
  borderline west={2.5pt}{0pt}{RuleGrey},
  colback=white, breakable, enhanced jigsaw,
  left=11pt,right=8pt,top=4pt,bottom=4pt,
  before skip=-4pt, after skip=12pt}
\newenvironment{proofomit}[1][MTH2025]{%
  \begin{omitbox}\textit{\color{RuleGrey}Proof omitted (#1).}\;\ignorespaces}{\end{omitbox}}
\makeatother

% ===================================================================
%  Math macros
% ===================================================================
\newcommand{\R}{\mathbb{R}}
\newcommand{\C}{\mathbb{C}}
\newcommand{\Q}{\mathbb{Q}}
\newcommand{\Z}{\mathbb{Z}}
\newcommand{\N}{\mathbb{N}}
\newcommand{\F}{\mathbb{F}}
\renewcommand{\vec}[1]{\mathbf{#1}}
\newcommand{\vu}{\mathbf{u}}
\newcommand{\vv}{\mathbf{v}}
\newcommand{\vw}{\mathbf{w}}
\newcommand{\vx}{\mathbf{x}}
\newcommand{\vy}{\mathbf{y}}
\newcommand{\vb}{\mathbf{b}}
\newcommand{\vc}{\mathbf{c}}
\newcommand{\vz}{\mathbf{z}}
\newcommand{\vze}{\mathbf{0}}
\newcommand{\ve}{\mathbf{e}}
\newcommand{\vp}{\mathbf{p}}
\newcommand{\veci}{\mathbf{i}}
\newcommand{\trans}{^{\mathsf{T}}}
\newcommand{\inv}{^{-1}}
\newcommand{\spn}{\operatorname{span}}
\newcommand{\rank}{\operatorname{rank}}
\newcommand{\nullity}{\operatorname{nullity}}
\newcommand{\tr}{\operatorname{tr}}
\newcommand{\rowsp}{\operatorname{row}}
\newcommand{\colsp}{\operatorname{col}}
\newcommand{\ran}{\operatorname{ran}}
\newcommand{\coker}{\operatorname{coker}}
\newcommand{\proj}{\operatorname{proj}}
\newcommand{\dist}{\operatorname{dist}}
\renewcommand{\dim}{\operatorname{dim}}
\newcommand{\diag}{\operatorname{diag}}
\newcommand{\Mat}{\mathrm{M}}
\newcommand{\inner}[2]{\langle #1, #2 \rangle}
\newcommand{\norm}[1]{\lVert #1 \rVert}
\newcommand{\abs}[1]{\lvert #1 \rvert}
\DeclareMathOperator{\REFop}{REF}
\DeclareMathOperator{\RREFop}{RREF}
\newcommand{\REF}{\ensuremath{\REFop}}
\newcommand{\RREF}{\ensuremath{\RREFop}}

% Cross-reference shorthand: the argument IS the seminar-notes number,
% so we print it literally and hyperlink to the environment's target.
\newcommand{\Thm}[1]{\hyperlink{thm:#1}{Theorem~#1}}
\newcommand{\Lem}[1]{\hyperlink{lem:#1}{Lemma~#1}}
\newcommand{\Cors}[1]{\hyperlink{cor:#1}{Corollary~#1}}
\newcommand{\Props}[1]{\hyperlink{prop:#1}{Proposition~#1}}
\newcommand{\Defn}[1]{\hyperlink{def:#1}{Definition~#1}}
\newcommand{\Alg}[1]{\hyperlink{alg:#1}{Algorithm~#1}}
\newcommand{\Rmk}[1]{\hyperlink{rmk:#1}{Remark~#1}}
\newcommand{\Egref}[1]{\hyperlink{eg:#1}{Example~#1}}
\newcommand{\axiom}[1]{\textsf{#1}}
